% Vietnamese translation for OI Public Library
% Source-File: IOI中国国家候选队论文/2005/魏冉/魏冉.doc
\documentclass[11pt]{article}
\usepackage{oipl-vn}
\usepackage{tikz}
\usetikzlibrary{arrows.meta,positioning}

\newcommand{\Oh}{\mathcal{O}}
\newcommand{\Prob}{\mathop{\mathrm{Pr}}}

\title{Làm cho hiệu suất thuật toán ``nhảy lên'': thao tác và ứng dụng của skip list}
\author{Wei Ran\\Trường Trung học số 2 trực thuộc Đại học Sư phạm Hoa Đông, Thượng Hải}
\date{2005}

\begin{document}
\maketitle

\origtitle{Làm cho hiệu suất thuật toán ``nhảy lên'': thao tác và ứng dụng của skip list}
\sourcepath{IOI China candidate team papers / 2005 / Wei Ran.doc}

\renewcommand{\abstractname}{Tóm tắt}
\begin{abstract}
Bài viết gồm ba phần chính. Phần đầu giới thiệu skip list từ chức năng, hiệu
suất và cấu trúc trực quan, để người đọc có hình dung ban đầu về cấu trúc dữ
liệu này. Phần thứ hai trình bày ba thao tác cơ bản: tìm kiếm, chèn và xóa,
đồng thời phân tích độ phức tạp thời gian và bộ nhớ của chúng. Phần thứ ba
giới thiệu ứng dụng của skip list, dùng kết quả thử nghiệm để so sánh skip
list với một số cấu trúc dữ liệu liên quan, từ đó nêu ưu nhược điểm của từng
cách tiếp cận. Cuối cùng, bài viết rút ra vài nhận xét về giá trị thực hành
của skip list.
\end{abstract}

\noindent\textbf{Từ khóa:} skip list; hiệu suất; xác suất; ngẫu nhiên hóa.

\section{Tổng quan và cấu trúc}

Cây nhị phân là một cấu trúc dữ liệu quen thuộc, hỗ trợ các thao tác như tìm
kiếm, chèn và xóa. Điểm yếu chí mạng của cây nhị phân tìm kiếm thông thường
là khi dữ liệu không đủ ngẫu nhiên, hình dạng cây có thể mất cân bằng, kéo
theo hiệu suất thuật toán giảm mạnh.

Skip list là một cấu trúc dữ liệu ra đời năm 1987. Với các thao tác tìm kiếm,
chèn và xóa, thời gian kỳ vọng đều là \(\Oh(\log n)\), gần như có thể thay thế
các cây cân bằng trong nhiều tình huống. Điểm quan trọng nhất là độ phức tạp
cài đặt của skip list thấp hơn nhiều so với AVL tree, red-black tree và các
cấu trúc cùng loại. Vì vậy nó có lợi thế rõ rệt cả về khả năng hiểu lẫn khả
năng phổ biến.

Một skip list gồm nhiều danh sách liên kết
\[
  S_0,S_1,S_2,\ldots,S_h
\]
và thỏa ba điều kiện:
\begin{itemize}
  \item Mỗi danh sách đều chứa hai phần tử đặc biệt \(-\infty\) và \(+\infty\).
  \item \(S_0\) chứa tất cả phần tử, và các phần tử trong mỗi danh sách được
  sắp xếp tăng dần.
  \item Với mọi \(i>0\), tập phần tử của \(S_i\) là tập con của tập phần tử
  trong \(S_{i-1}\).
\end{itemize}

Hình trong bản gốc minh họa một skip list có bốn tầng \(S_0\) tới \(S_3\).
Các phần tử xuất hiện ở tầng cao cũng xuất hiện ở mọi tầng thấp hơn, tạo thành
những ``cột'' thẳng đứng. Khi tìm kiếm, ta có thể đi ngang trên tầng cao để bỏ
qua nhiều phần tử, rồi đi xuống dần khi cần tinh chỉnh vị trí.

\begin{figure}[ht]
\centering
\begin{tikzpicture}[
  >=Latex,
  node/.style={draw, minimum width=8mm, minimum height=6mm, inner sep=1pt, font=\small},
  level/.style={font=\small},
  x=1.05cm, y=0.62cm
]
\foreach \y/\name in {0/S_0,1/S_1,2/S_2,3/S_3} {
  \node[level] at (-1,\y) {\(\name\)};
}
\foreach \x/\v in {0/{-\infty},1/10,2/20,3/35,4/40,5/53,6/70,7/{+\infty}} {
  \node[node] (b\x) at (\x,0) {\(\v\)};
}
\foreach \x/\v in {0/{-\infty},2/20,4/40,6/70,7/{+\infty}} {
  \node[node] (c\x) at (\x,1) {\(\v\)};
}
\foreach \x/\v in {0/{-\infty},4/40,7/{+\infty}} {
  \node[node] (d\x) at (\x,2) {\(\v\)};
}
\foreach \x/\v in {0/{-\infty},4/40,7/{+\infty}} {
  \node[node] (e\x) at (\x,3) {\(\v\)};
}
\foreach \a/\b in {b0/b1,b1/b2,b2/b3,b3/b4,b4/b5,b5/b6,b6/b7,
                  c0/c2,c2/c4,c4/c6,c6/c7,
                  d0/d4,d4/d7,e0/e4,e4/e7} {
  \draw[->] (\a) -- (\b);
}
\foreach \x in {0,4,7} {
  \draw[densely dotted] (b\x) -- (e\x);
}
\foreach \x in {2,6} {
  \draw[densely dotted] (b\x) -- (c\x);
}
\node[below=6mm of b4, align=center, font=\small]
  {Nút ở tầng cao tạo thành đường tắt, còn \(S_0\) vẫn giữ toàn bộ thứ tự.};
\end{tikzpicture}
\caption{Ví dụ skip list bốn tầng: các cột cao hơn xuất hiện thưa dần theo xác suất.}
\end{figure}

\section{Các thao tác cơ bản}

\subsection{Tìm kiếm}

Mục tiêu là tìm một phần tử \(x\) trong skip list. Quy trình như sau:
\begin{enumerate}
  \item Bắt đầu từ đầu danh sách ở tầng cao nhất \(S_h\).
  \item Giả sử vị trí hiện tại là \(p\), nút bên phải mà \(p\) trỏ tới là
  \(q\), và giá trị của \(q\) là \(y\). So sánh \(y\) với \(x\):
  \begin{itemize}
    \item nếu \(x=y\), tìm kiếm thành công;
    \item nếu \(x>y\), di chuyển sang phải tới \(q\);
    \item nếu \(x<y\), di chuyển xuống một tầng.
  \end{itemize}
  \item Nếu đã ở tầng đáy \(S_0\) mà vẫn cần đi xuống, tìm kiếm thất bại.
\end{enumerate}

Bản gốc dùng ví dụ tìm phần tử \(53\): đường đi bắt đầu từ tầng cao nhất, đi
ngang qua các nút nhỏ hơn \(53\), rồi hạ tầng khi nút tiếp theo đã vượt mục
tiêu. Đây chính là cơ chế ``nhảy'' làm skip list nhanh hơn danh sách liên kết
đơn.

\begin{figure}[ht]
\centering
\begin{tikzpicture}[
  >=Latex,
  node/.style={draw, minimum width=8mm, minimum height=6mm, inner sep=1pt, font=\small},
  x=1.05cm, y=0.62cm
]
\foreach \y/\name in {0/S_0,1/S_1,2/S_2,3/S_3} {
  \node[font=\small] at (-1,\y) {\(\name\)};
}
\foreach \x/\v in {0/{-\infty},1/10,2/20,3/35,4/40,5/53,6/70,7/{+\infty}} {
  \node[node] (b\x) at (\x,0) {\(\v\)};
}
\foreach \x/\v in {0/{-\infty},2/20,4/40,6/70,7/{+\infty}} {
  \node[node] (c\x) at (\x,1) {\(\v\)};
}
\foreach \x/\v in {0/{-\infty},4/40,7/{+\infty}} {
  \node[node] (d\x) at (\x,2) {\(\v\)};
}
\foreach \x/\v in {0/{-\infty},4/40,7/{+\infty}} {
  \node[node] (e\x) at (\x,3) {\(\v\)};
}
\foreach \a/\b in {b0/b1,b1/b2,b2/b3,b3/b4,b4/b5,b5/b6,b6/b7,
                  c0/c2,c2/c4,c4/c6,c6/c7,
                  d0/d4,d4/d7,e0/e4,e4/e7} {
  \draw[->, gray!55] (\a) -- (\b);
}
\foreach \x in {0,4,7} {
  \draw[densely dotted, gray!60] (b\x) -- (e\x);
}
\foreach \x in {2,6} {
  \draw[densely dotted, gray!60] (b\x) -- (c\x);
}
\draw[->, very thick, red!70!black] (e0) -- (e4);
\draw[->, very thick, red!70!black] (e4) -- (d4);
\draw[->, very thick, red!70!black] (d4) -- (c4);
\draw[->, very thick, red!70!black] (c4) -- (b4);
\draw[->, very thick, red!70!black] (b4) -- (b5);
\node[above=5mm of e4, font=\small, red!70!black] {đường tìm \(53\)};
\end{tikzpicture}
\caption{Tìm \(53\): đi ngang khi nút kế tiếp chưa vượt mục tiêu, và đi xuống khi cần tinh chỉnh.}
\end{figure}

\subsection{Chèn}

Chèn một phần tử \(x\) tương đương với việc thêm một cột liên tiếp bắt đầu từ
tầng \(S_0\). Có hai tham số cần xác định: vị trí của cột và chiều cao của cột.

Vị trí được xác định bằng thao tác tìm kiếm: tìm phần tử lớn nhất \(y<x\).
Do mọi danh sách đều tăng dần, \(x\) phải được chèn ngay sau \(y\).

Chiều cao của cột quan trọng hơn và khó xác định hơn. Nếu chọn không tốt,
hiệu suất của cấu trúc có thể thay đổi lớn. Để sau khi chèn dữ liệu, các thao
tác vẫn giữ thời gian kỳ vọng \(\Oh(\log n)\), skip list dùng một quyết định
ngẫu nhiên hóa. Với một hằng số \(p\in(0,1)\), mô-đun quyết định có dạng:
\[
  r\leftarrow \mathrm{random}(0,1),\qquad
  \begin{cases}
    \text{thực hiện phương án A}, & r<p,\\
    \text{thực hiện phương án B}, & r\ge p.
  \end{cases}
\]

Ban đầu chiều cao cột là \(1\). Mỗi lần mô-đun trả về A, tăng chiều cao thêm
1 và tiếp tục gieo. Khi mô-đun trả về B ở lần thứ \(i\), dừng lại và chèn một
cột cao \(i\). Khi đó:
\[
  \Prob(\text{chiều cao cột}\ge k)=p^{k-1}.
\]
Nếu chiều cao mới lớn hơn chiều cao hiện tại \(h\) của skip list, cần thêm các
danh sách mới phía trên để cấu trúc vẫn thỏa các điều kiện đã nêu.

Bản gốc minh họa thao tác chèn \(40\): trước hết tìm phần tử lớn nhất nhỏ hơn
\(40\), rồi chèn một cột cao \(4\). Nếu chiều cao này vượt chiều cao cũ, một
tầng mới được tạo thêm.

\begin{figure}[ht]
\centering
\begin{tikzpicture}[
  >=Latex,
  node/.style={draw, minimum width=8mm, minimum height=6mm, inner sep=1pt, font=\small},
  newnode/.style={node, fill=gray!15, very thick},
  x=1.05cm, y=0.62cm
]
\foreach \y/\name in {0/S_0,1/S_1,2/S_2,3/S_3} {
  \node[font=\small] at (-1,\y) {\(\name\)};
}
\foreach \x/\v in {0/{-\infty},1/10,2/20,3/35,5/53,6/70,7/{+\infty}} {
  \node[node] (b\x) at (\x,0) {\(\v\)};
}
\foreach \x/\v in {0/{-\infty},2/20,6/70,7/{+\infty}} {
  \node[node] (c\x) at (\x,1) {\(\v\)};
}
\foreach \x/\v in {0/{-\infty},7/{+\infty}} {
  \node[node] (d\x) at (\x,2) {\(\v\)};
  \node[node] (e\x) at (\x,3) {\(\v\)};
}
\foreach \y/\prefix in {0/b,1/c,2/d,3/e} {
  \node[newnode] (n\y) at (4,\y) {\(40\)};
}
\foreach \a/\b in {b0/b1,b1/b2,b2/b3,b3/n0,n0/b5,b5/b6,b6/b7,
                  c0/c2,c2/n1,n1/c6,c6/c7,
                  d0/n2,n2/d7,e0/n3,n3/e7} {
  \draw[->] (\a) -- (\b);
}
\draw[densely dotted, very thick] (n0) -- (n3);
\node[below=6mm of n0, align=center, font=\small]
  {Chiều cao ngẫu nhiên bằng \(4\), nên nút \(40\) được nối vào bốn tầng.};
\end{tikzpicture}
\caption{Chèn \(40\): sau khi tìm vị trí, tạo một cột có chiều cao do mô-đun ngẫu nhiên quyết định.}
\end{figure}

\subsection{Xóa}

Xóa phần tử \(x\) gồm ba bước:
\begin{enumerate}
  \item Tìm vị trí của \(x\) trong skip list; nếu không tồn tại thì kết thúc.
  \item Xóa toàn bộ cột chứa \(x\) khỏi các danh sách.
  \item Xóa các tầng trên cùng đã trở thành rỗng, tức chỉ còn hai nút biên
  \(-\infty\) và \(+\infty\).
\end{enumerate}

\begin{figure}[ht]
\centering
\begin{tikzpicture}[
  >=Latex,
  node/.style={draw, minimum width=8mm, minimum height=6mm, inner sep=1pt, font=\small},
  removed/.style={node, dashed, gray!70},
  x=1.05cm, y=0.62cm
]
\foreach \y/\name in {0/S_0,1/S_1} {
  \node[font=\small] at (-1,\y) {\(\name\)};
}
\foreach \x/\v in {0/{-\infty},1/10,2/20,3/35,4/40,5/53,6/70,7/{+\infty}} {
  \node[node] (b\x) at (\x,0) {\(\v\)};
}
\foreach \x/\v in {0/{-\infty},2/20,4/40,5/53,6/70,7/{+\infty}} {
  \node[node] (c\x) at (\x,1) {\(\v\)};
}
\node[removed] at (5,0) {\(53\)};
\node[removed] at (5,1) {\(53\)};
\foreach \a/\b in {b0/b1,b1/b2,b2/b3,b3/b4,b4/b6,b6/b7,
                  c0/c2,c2/c4,c4/c6,c6/c7} {
  \draw[->] (\a) -- (\b);
}
\draw[densely dotted, gray!60] (b5) -- (c5);
\draw[->, very thick, red!70!black] (b4) to[bend right=18] (b6);
\draw[->, very thick, red!70!black] (c4) to[bend left=18] (c6);
\node[below=7mm of b5, align=center, font=\small]
  {Xóa toàn bộ cột \(53\), rồi nối lại các tiền nhiệm với hậu nhiệm trên từng tầng.};
\end{tikzpicture}
\caption{Xóa một phần tử khỏi skip list là xóa cả cột của nó và vá lại các liên kết ngang.}
\end{figure}

\subsection{Tìm kiếm có nhớ}

Tìm kiếm có nhớ, thường gọi là \emph{search with fingers}, là cách tận dụng
thông tin từ lần tìm kiếm trước để hỗ trợ lần tìm kiếm kế tiếp. Nếu hai phần
tử được tìm liên tiếp nằm gần nhau trong thứ tự của skip list, kỹ thuật này
có thể giảm độ phức tạp kỳ vọng xuống \(\Oh(\log k)\), trong đó \(k\) là
khoảng cách thứ tự giữa hai phần tử ấy.

Giả sử lần trước tìm phần tử \(i\), lần này cần tìm phần tử \(j\). Ta dùng
mảng \(\textit{update}\) để ghi lại, trên mỗi tầng, vị trí bên phải nhất mà
đường tìm kiếm đã nhảy tới khi tìm \(i\). Khi tìm \(j\):
\begin{itemize}
  \item Nếu \(i\le j\), bắt đầu từ tầng \(S_0\), duyệt ngược lên trong mảng
  \(\textit{update}\) cho tới khi gặp một vị trí có nút bên phải mang giá trị
  ít nhất \(j\), rồi từ đó chạy quy trình tìm kiếm thông thường.
  \item Nếu \(i>j\), cũng duyệt từ tầng \(S_0\) lên trên, tìm vị trí thích hợp
  mà từ đó phần tử bên trái hoặc hiện tại không vượt quá \(j\), rồi tiếp tục
  tìm kiếm thông thường.
\end{itemize}

Kỹ thuật này đặc biệt hiệu quả với dữ liệu có tương quan mạnh giữa các lần
truy vấn liên tiếp. Với dữ liệu gần như ngẫu nhiên, chi phí duy trì và quét
thông tin nhớ có thể không đem lại lợi ích.

\section{Phân tích độ phức tạp}

Các đại lượng quan trọng của skip list đều được hiểu theo kỳ vọng:
\[
\begin{array}{ll}
\text{Bộ nhớ} & \Oh(n),\\
\text{Chiều cao} & \Oh(\log n),\\
\text{Tìm kiếm} & \Oh(\log n),\\
\text{Chèn} & \Oh(\log n),\\
\text{Xóa} & \Oh(\log n).
\end{array}
\]
Lý do cần nói ``kỳ vọng'' là vì phân tích của skip list dựa trên xác suất. Có
thể xuất hiện trường hợp xấu, nhưng xác suất ấy rất nhỏ.

\subsection{Bộ nhớ}

Giả sử có \(n\) phần tử. Theo công thức chiều cao, xác suất một phần tử xuất
hiện ở tầng \(S_i\) là \(p^{i-1}\). Vì vậy số phần tử kỳ vọng trên tầng \(i\)
là \(n p^{i-1}\), và tổng số nút kỳ vọng là
\[
  \sum_{i\ge 1} n p^{i-1}=\frac{n}{1-p}.
\]
Khi \(p<1/2\), tổng này nhỏ hơn \(2n\). Do đó bộ nhớ kỳ vọng là \(\Oh(n)\).

\subsection{Chiều cao}

Tương tự, ở tầng cao \(1+3\log_{1/p} n\), số phần tử kỳ vọng xấp xỉ
\[
  n p^{3\log_{1/p}n}=\frac{1}{n^2}.
\]
Khi \(n\) lớn, giá trị này gần \(0\), nên xác suất skip list còn phần tử ở cao
hơn mức ấy là rất nhỏ. Vì vậy chiều cao kỳ vọng thuộc cấp \(\Oh(\log n)\).

\subsection{Tìm kiếm}

Ta dùng phân tích ngược. Giả sử đang đứng ở nút đích và muốn đi ngược về nút
bắt đầu ở góc trái trên. Độ dài đường đi ngược này tương ứng với số bước của
tìm kiếm.

Xét một nút ở tầng \(i\). Nếu cột của nó đúng cao \(i\), sự kiện này có xác
suất \(1-p\), và đường đi ngược phải đến từ bên trái. Nếu cột cao hơn \(i\),
sự kiện này có xác suất \(p\), và đường đi ngược đến từ phía trên. Gọi
\(C(k)\) là số bước kỳ vọng để đi ngược lên \(k\) tầng, kể cả các bước ngang.
Khi đó:
\[
\begin{aligned}
  C(0)&=0,\\
  C(k)&=(1-p)(1+C(k))+p(1+C(k-1)).
\end{aligned}
\]
Suy ra
\[
  C(k)=\frac{1}{p}+C(k-1)=\frac{k}{p}.
\]
Vì chiều cao \(k\) ở cấp \(\Oh(\log n)\), thời gian tìm kiếm kỳ vọng cũng là
\(\Oh(\log n)\). Với tìm kiếm có nhớ, lập luận tương tự cho kết quả
\(\Oh(\log k)\), trong đó \(k\) là khoảng cách giữa hai phần tử được tìm liên
tiếp.

\subsection{Chèn và xóa}

Chèn và xóa gồm hai phần: tìm vị trí và cập nhật liên kết. Phần tìm kiếm có
độ phức tạp \(\Oh(\log n)\), còn phần cập nhật tỷ lệ với chiều cao của skip
list, cũng là \(\Oh(\log n)\). Do đó cả chèn lẫn xóa đều có thời gian kỳ vọng
\(\Oh(\log n)\).

\subsection{Kết quả thử nghiệm}

Bản gốc chạy \(10^6\) thao tác ngẫu nhiên để so sánh các giá trị \(p\):
\[
\begin{array}{lrrrrr}
p & \text{thời gian} & \text{cao TB} & \text{số nút} &
\text{nhảy tìm} & \text{nhảy chèn/xóa}\\
\hline
2/3 & 0.0024690\,\mathrm{ms} & 3.0049 & 123339 & 39.878 & 41.604/41.566\\
1/2 & 0.0020180\,\mathrm{ms} & 1.9956 & 68327 & 27.807 & 29.947/29.072\\
1/e & 0.0019870\,\mathrm{ms} & 1.5844 & 57027 & 27.332 & 28.238/28.452\\
1/4 & 0.0021720\,\mathrm{ms} & 1.3304 & 47828 & 28.726 & 29.472/29.664\\
1/8 & 0.0026880\,\mathrm{ms} & 1.1443 & 44203 & 35.147 & 35.821/36.007
\end{array}
\]
Kết quả cho thấy \(p=1/2\) và \(p=1/e\) có hiệu suất thời gian cao. Nếu bộ nhớ
là ràng buộc chặt, có thể chọn \(p\) nhỏ hơn, chẳng hạn \(1/4\).

Với tìm kiếm có nhớ, bản gốc ghi nhận:
\[
\begin{array}{llrrrr}
p & \text{dữ liệu} & \text{không nhớ} & \text{có nhớ} &
\text{nhảy không nhớ} & \text{nhảy có nhớ}\\
\hline
0.5 & \text{ngẫu nhiên} & 0.0020150 & 0.0020790 & 23.262 & 26.509\\
0.5 & \text{có tương quan} & 0.0008440 & 0.0006880 & 26.157 & 4.932
\end{array}
\]
Đơn vị thời gian là mili giây trên mỗi thao tác. Khi các khóa được tìm liên
tiếp gần nhau, tìm kiếm có nhớ đem lại lợi ích rõ rệt. Khi dữ liệu ngẫu nhiên
mạnh, kỹ thuật này không những không nhanh hơn mà còn có thể trở thành nút
thắt do số lần nhảy tăng.

\section{Ứng dụng của skip list}

Hiệu suất cao và độ phức tạp cài đặt thấp làm skip list hữu ích trong nhiều
bài toán thực tế, đặc biệt khi thời gian lập trình rất hạn chế. Bài viết dùng
hai bài thi để minh họa.

\subsection{NOI 2004 Day 1: Cashier}

Bài Cashier yêu cầu duy trì hồ sơ lương nhân viên dưới bốn thao tác: thêm nhân
viên mới, tăng toàn bộ lương, giảm toàn bộ lương và hỏi mức lương lớn thứ
\(k\). Khi lương bị giảm xuống dưới ngưỡng hợp đồng, nhân viên rời công ty và
phải bị xóa khỏi cấu trúc.

Gọi \(R\) là miền giá trị lương và \(N\) là số nhân viên. Ba hướng giải được
so sánh như sau:
\[
\begin{array}{lllll}
\text{Cấu trúc} & I & A & S & F\\
\hline
\text{segment tree theo lương} & \Oh(\log R) & \Oh(1) & \Oh(\log R) & \Oh(\log R)\\
\text{splay tree theo lương} & \Oh(\log N) & \Oh(1) & \Oh(\log N) & \Oh(\log N)\\
\text{skip list} & \Oh(\log N) & \Oh(1) & \Oh(\log N) & \Oh(\log N)
\end{array}
\]

Kết quả chạy thực tế trong bản gốc, đơn vị giây:
\[
\begin{array}{lrrrrrrrrrr}
\text{Cấu trúc} & T1&T2&T3&T4&T5&T6&T7&T8&T9&T10\\
\hline
\text{segment tree} & 0.000&0.000&0.000&0.031&0.062&0.094&0.109&0.203&0.265&0.250\\
\text{splay tree} & 0.000&0.000&0.016&0.062&0.047&0.125&0.141&0.360&0.453&0.422\\
\text{skip list} & 0.000&0.000&0.000&0.047&0.062&0.109&0.156&0.368&0.438&0.375
\end{array}
\]

Segment tree tỏ ra nhanh nhất trong dữ liệu này, nhưng nó dựa trực tiếp vào
miền khóa. Ở bài gốc, \(R\) chỉ khoảng \(10^5\), nên dùng segment tree còn
hợp lý. Nếu miền khóa lớn hơn nhiều, hoặc khóa không phải số nguyên, segment
tree trở nên khó áp dụng. Khi đó các cấu trúc dựa trên phần tử như splay tree
và skip list thích hợp hơn. Kết quả thực nghiệm cũng cho thấy skip list không
kém splay tree, trong khi cài đặt đơn giản hơn.

\subsection{HNOI 2004 Day 1: Pet Adoption}

Bài Pet Adoption duy trì một tập thú nuôi hoặc một tập người nhận nuôi. Mỗi
lần có đối tượng mới đến, nếu phía còn lại không rỗng thì ghép với đối tượng
có trị đặc trưng gần nhất; nếu hòa, chọn trị nhỏ hơn. Miền trị đặc trưng lên
tới \(2^{31}\).

Điểm khác lớn nhất so với Cashier là miền khóa quá rộng. Segment tree gặp áp
lực bộ nhớ lớn; dù dùng cách cấp phát nút động, nếu các khóa phân tán thì độ
phức tạp bộ nhớ có thể tiến gần \(\Oh(N\log N)\). Nếu mở rộng bài toán để
khóa là số thực, segment tree gần như mất tác dụng.

Skip list phù hợp tự nhiên với bài này. Hầu như không cần sửa thuật toán
chuẩn, có thể viết nhanh nếu đã quen, hiệu suất tốt, và quan trọng hơn là nó
không phụ thuộc chặt vào kiểu khóa. Tính mở rộng này là một ưu điểm lớn của
skip list trong thi đấu.

\section{Kết luận}

Skip list là một cấu trúc dữ liệu mới mẻ so với nhiều cây cân bằng kinh điển,
nhưng có hiệu suất cao và độ phức tạp cài đặt thấp. So với splay tree, vốn
cũng nổi tiếng vì tương đối dễ viết, skip list không hề kém về hiệu suất, và
trong một số thử nghiệm còn tốt hơn.

Khi suy nghĩ về một lớp bài toán, ta thường vô thức bị giới hạn trong một
phạm vi quen thuộc. Với các bài toán liên quan tới cây cân bằng, nhiều cấu
trúc mạnh như red-black tree hay AVL tree đã được tạo ra, nhưng chúng vẫn nằm
trong khuôn khổ ``cây'' và thường có cài đặt phức tạp. Skip list cho thấy một
cấu trúc không liên quan trực tiếp đến cây vẫn có thể thực hiện chức năng
tương tự.

Cái tên ``skip list'' vì thế không chỉ mô tả hình dạng dữ liệu, mà còn gợi ra
một cách nghĩ: khi gặp bế tắc, hãy quay về bản chất của vấn đề và thử nhảy ra
khỏi khuôn mẫu quen thuộc.

\section{Phụ lục}

\subsection{Vì sao p bằng một trên e thường hiệu quả?}

Từ phân tích độ phức tạp, thời gian phụ thuộc vào số lần nhảy, xấp xỉ
\[
  \frac{k}{p},
\]
trong đó \(k\) là chiều cao skip list. Mà
\[
  k=\Theta(\log_{1/p} n).
\]
Bỏ qua hằng số chỉ phụ thuộc vào \(n\), cần tối ưu
\[
  f(p)=\frac{1}{p}\log_{1/p} n.
\]
Đặt \(x=1/p\). Khi đó, về phần phụ thuộc vào \(x\),
\[
  g(x)=\frac{x}{\ln x}.
\]
Đạo hàm cho
\[
  g'(x)=\frac{\ln x-1}{(\ln x)^2}.
\]
Do đó \(g'(x)=0\) khi \(x=e\), tức \(p=1/e\). Đây là lý do thực nghiệm thường
cho thấy \(p=1/e\) đạt hiệu suất thời gian rất tốt.

\subsection{Ghi chú về phụ lục chương trình và hình vẽ}

Bản gốc có bảng so sánh lấy từ bài ``Skip Lists: A Probabilistic Alternative
to Balanced Trees'' của William Pugh, một chương trình Pascal cài đặt skip
list, cùng các hình minh họa thao tác tìm kiếm, chèn và xóa. Các hình thao tác
chính đã được vẽ lại ở các phần trên. Mã Pascal không phải nội dung cần dịch,
nên không được chép lại đầy đủ.

\subsection{Hai đề bài được dùng làm ví dụ}

Cashier yêu cầu xử lý tối đa \(100000\) lệnh thêm nhân viên và \(100000\) truy
vấn hạng lương. Có bốn loại lệnh: thêm hồ sơ với lương ban đầu, tăng lương
toàn công ty, giảm lương toàn công ty và hỏi lương lớn thứ \(k\). Cần in đáp
án cho từng truy vấn và tổng số nhân viên đã rời công ty.

Pet Adoption yêu cầu xử lý tối đa \(80000\) sự kiện. Mỗi sự kiện là một thú
nuôi hoặc một người nhận nuôi với trị đặc trưng dương nhỏ hơn \(2^{31}\). Nếu
phía đối ứng đang chờ, ghép với trị gần nhất và cộng độ bất mãn
\(|a-b|\); nếu không, đưa đối tượng mới vào tập chờ. Cần in tổng độ bất mãn
lấy modulo \(1000000\).

\begin{thebibliography}{9}
\bibitem{pugh-skip-list}
William Pugh. \emph{Skip Lists: A Probabilistic Alternative to Balanced Trees}.
\bibitem{pugh-cookbook}
William Pugh. \emph{A Skip List Cookbook}.
\end{thebibliography}

\end{document}
